\section{Differentiaaligeometriaa}
Käyrien muodostamien linkkien tutkimiseen tarvitaan (sileiden) monistojen teoriaa. Monistot ovat geometrisia rakenteita, jotka muistuttavat paikallisesti euklidista avaruutta. Ne mahdollistavat \emph{avaruuden rakenteen} käsitteen. Algebrallisten solmujen kannalta tämä on olennaista, sillä solmujen tutkimista helpottaa niiden upottaminen monistoon. Tässä työssä noudatetaan lähteen \citep{Lee2012} merkintätapoja.
\begin{huomautus}
    Tässä osiossa, kuten kirjallisuudessa yleisesti, käsitellään $\R$-monistoja, vaikka työ muuten käsittelee $\C$-käyriä. Monistojen teoria sellaisenaan ei kuitenkaan välitä skalaareiden reaalisuudesta, ja isomorfia $\C\cong\R^2$ sallii teorian siirtämisen kompleksiluvuille.
\end{huomautus}

\subsection{Sileät monistot ja Seifert-monisto}
Sileiden monistojen määrittely on jotakuinkin teknistä niiden ymmärrettävyyteen verrattuna. Kärsimätön lukija voi ohittaa osion, ja ajatella monistoja $n$-ulotteisina pintoina, jotka $n+1$-ulotteiseen avaruuteen upotettuna näyttävät paikallisesti $n$-ulotteisilta tasoilta, mutta saattavat globaalisti olla vääntyneitä.

Sileä monisto on sellainen monisto, joka ei ole ``rutussa'', jolloin sen pinnalla voi integroida ja derivoida käyriä. Esimerkiksi maapallo on sileä 2-ulotteinen monisto.
\begin{maaritelma}
    Olkoon $M$ topologinen avaruus. $M$ on \emph{$n$-ulotteinen topologinen monisto}, jos
    \begin{itemize}
        \item $M$ on Hausdorff-avaruus
        \item $M$ on $N_2$-avaruus
        \item $M$ on paikallisesti euklidisesti $n$-ulotteinen: kaikilla $p\in M$ on avoin ympäristö $U\subset M$, joka on homeomorfinen johonkin $\R^n$:n avoimeen osajoukkoon $V\subset\R^n$.
    \end{itemize}
\end{maaritelma}
\begin{huomautus}
    Yleisesti kirjoitetaan ``olkoon $M^n$ monisto'', jolla tarkoitetaan, että $M$ on $n$-ulotteinen monisto.
\end{huomautus}
\begin{esimerkki}
    Esimerkiksi Maapallo on 2-ulotteinen monisto, koska paikallisesti Maapallon pinta muistuttaa tasoa. Maapallo on erityistapaus algebrallisesta pinnasta, koska sen polynomiyhtälö $x^2+y^2+z^2=R^2$.
\end{esimerkki}
\begin{maaritelma}
    Olkoon $M^n$ monisto. $M$:n \emph{koordinaattikartta} on pari $(U,\varphi)$, jossa $U\subset M$ on avoin ja $\varphi:U\to\hat U$ on homeomorfismi avoimeen joukkoon $\hat U\subset R^n$. Koordinaattikartan \emph{koordinaattikuvaus} on $\varphi$.
\end{maaritelma}
\begin{esimerkki}
    Yleisesti algebrallinen hyperpinta $f(x_1,x_2,\dots,x_{n+1})=0$ määrittää $n$-ulotteisen moniston $\R^{n+1}$:ssä.
\end{esimerkki}
\begin{maaritelma}
    Jos $U\subset\R^n,V\subset\R^m$, funktio $f:U\to V$ on \emph{sileä}, jos sen jokaisella komponenttifunktiolla on jatkuvat asteen $n$ osittaisderivaatat kaikilla $n\in\N$.
\end{maaritelma}
\begin{maaritelma}
    Kuvaus $f:A\to B$ on \emph{diffeomorfismi}, jos $f$ on bijektio ja $f$ ja $f^{-1}$ ovat sileitä.
\end{maaritelma}
\begin{maaritelma}
    Moniston $M$ \emph{atlas} $\mathcal{A}$ on kokoelma karttoja $\{(U_i,\varphi_i)\}_{i\in I}$ siten, että karttojen lähtöjoukot peittävät moniston $M$: kaikille $p\in M$ on olemassa $i\in I$, jolle $p\in U_i$. Jos koordinaattikartat ovat pareittain sileästi yhteensopivia, eli
    \begin{equation}
        \varphi_j\circ\varphi_i^{-1}:\varphi_i(U_i\cap U_j)\to\varphi_j(U_i\cap U_j)
    \end{equation}
    tai $U_i\cap U_j=\emptyset$, niin $\mathcal{A}$ on \emph{sileä atlas}.
\end{maaritelma}
\begin{maaritelma}
    Monisto $M$ on \emph{sileä monisto}, jos sillä on sileä atlas.
\end{maaritelma}
Osion loppuhuipennuksena määritellään Seifert-monisto. 3-monistojen modernia teoriaa tuntemattomalle lukijalle määritelmä on käsittämätön, mutta käytännössä määritelmällä eritellään tutkimuksen kohteeksi mahdollisimman laaja monistojen joukko, joka käyttäytyisi mahdollisimman hyvin.
\begin{maaritelma}[Seifert-monisto ja Seifert-kuidutettu avaruus]
    \emph{Seifert-monisto} on 3-monisto, jonka voi pujontaoperaatiolla (määritelmä~\ref{maar:pujos}) rakentaa \emph{Seifert-kuidutetuista avaruuksista}, ts. avaruuksista, jotka voidaan esittää erillisten ympyröiden $\S^1$ yhdisteenä.
\end{maaritelma}
\begin{esimerkki}
    Sekä ontto että täytetty torus $\S^1\times\S^1$, $\S^1\times D^2$ ovat Seifert-kuidutettuja avaruuksia.

    $\S^3$ on Seifert-monisto, koska $\S^3=\T^2\sqcup\T^2$. Tarkalleen ottaen $\S^3$ saadaan pujomalla yhteen kaksi triviaalia solmua $(\S^3,\S^1)$. Siis $\S^3$ voidaan hajottaa kahteen joukkoon $\S^3-\S^1$, joille pätee $\S^3-\S^1\approx\interior\T^2$.
\end{esimerkki}

\subsection{Tangenttiavaruus}
Keskeinen osa sileiden monistojen motivaatiota on derivoimisen mahdollistaminen avaruuksissa, jotka ovat globaalisti kaarevia. Alunperin derivointi määriteltiin käyrille $\R^2$:ssa tangenttien avulla. Tässä tilanteessa käyrää voidaan ajatella sileänä $1$-monistona. Mitä jos käyrä ei olekaan upotettu $R^2$:een? Käyrälle ajatus voi tuntua naurettavalta, mutta moniulotteisemmissa monistoissa tilanne ei ole enää selvä. Esimerkiksi maailmankaikkeus on jonkinlainen monisto, mutta mihin avaruuteen se on upotettu?

On selvää, että tangentti ei sellaisenaan yleisty monistoille ja tarvitaan tarkempi määritelmä. Hyvän määritelmän ymmärtäisi myös moniston ``sisällä'' elävä olio. Tangenttiavaruuden määrittely alkaa derivaatan yleistyksestä derivaatiosta.
\begin{maaritelma}
    Olkoon $M$ sileä $n$-monisto ja $p\in M$. Olkoon $v:C^\infty(M)\to\R$ lineaarinen kuvaus. Jos se tyydyttää Leibniz-säännön
    \begin{equation}
        v(fg)=f(p)vg+g(p)vf,\quad\text{kaikille }f,g\in\C^\infty(M),
    \end{equation}
    niin $v$ on \emph{derivaatio} $p$:ssä.
\end{maaritelma}
\begin{maaritelma}
    Kaikkien derivaatioiden joukko $p$:ssä, merkitään $T_pM$, on $M$:n \emph{tangenttiavaruus} $p$:ssä. Jos $v\in T_pM$, niin $v$ on \emph{tangenttivektori} $p$:ssä.
\end{maaritelma}
Joissakin tapauksissa hyödyllinen on myös konsepti, joka niputtaa kaikki tangenttiavaruudet yhteen.
\begin{maaritelma}
    Moniston $M$ \emph{tangenttikimppu} on erillinen yhdiste
    \begin{equation}
        TM=\bigsqcup_{p\in M}T_pM.
    \end{equation}
\end{maaritelma}


\subsection{Vektoriavaruuden suunnistus}
Toisin kuin tutuissa vektoriavaruuksissa $\R^3\cong\text{span}\{e_1,e_2,e_3\}$, yleisessä vektoriavaruudessa $V$ ei ole päivänselvää tapaa määrittää ``positiivista'' tai ``negatiivista'' suuntaa.
\begin{esimerkki}
    Olkoon $V$ 2-ulotteinen vektoriavaruus, jonka kanta on $\{\sin\theta,\cos\theta\}$. Kumpi suunnistetuista kannoista $(\sin\theta,\cos\theta)$ ja $(\cos\theta,\sin\theta)$ on positiivinen?
\end{esimerkki}
Vektoriavaruuden tai moniston oikeaoppinen suunnistus voidaan kuitenkin kääntää muodollisiksi matemaattisiksi ilmaisuiksi \emph{differentiaalimuotojen} avulla. Aloitetaan kuitenkin hyvin määritellyistä suunnistuksista.
\begin{maaritelma}
    Olkoon $V$ reaalinen $n$-vektoriavaruus. Kaksi suunnistettua kantaa $(e_1,\dots,e_n)$ ja $(f_1,\dots,f_n)$ ovat \emph{yhtäpitävästi suunnistettuja}, jos niiden muuntomatriisille, jonka määrittää yhtälö
    \begin{equation}
        e_i=B_i^j f_j,
    \end{equation}
    pätee $\det(B_i^j)>0$.
\end{maaritelma}
Tämä muuntomatriisi määrittää ekvivalenssirelaation, joka jakaa kannat kahteen joukkoon. 
\begin{lemma}
    Vektoriavaruuden yhtäpitävä suunnistus on sen kantojen ekvivalenssirelaatio, jolla on tasan kaksi ekvivalenssiluokkaa.
\end{lemma}
\begin{proof}
    ???
\end{proof}
Matemaatikon tehtäväksi jää määritellä toinen näistä luokista ``positiiviseksi'' ja toinen ``negatiiviseksi''. Tämän jälkeen voidaan määritellä suunnistus vektoriavaruudessa.

\subsection{Suunnistus monistossa}
Moniston $M$ suunnistus määritellään sen tangenttiavaruuden suunnistuksen kautta.
\begin{maaritelma}
    Olkoon $M$ sileä $n$-monisto. $M$:n \emph{pisteittäinen suunnistus} on suunnistuksen valinta jokaiselle tangenttiavaruudelle $T_pM$, $p\in M$.
\end{maaritelma}
Sellaisenaan tämä ei ole järin hyödyllinen määritelmä, sillä suunnistus voi vaihdella pisteestä toiseen. Toivottavaa olisi etenkin, että moniston sileillä alimonistoilla (käyrillä), suunnistus muuttuisi sileästi alimoniston (käyrän) kaartuessa.
\begin{maaritelma}[Local frame, paikallinen kehys]
    ???
\end{maaritelma}
\begin{maaritelma}
    Olkoon $(e_i)$ tangenttikimpun $TM$ paikallinen kehys. Tällöin $(e_i)$ on \emph{(positiivisesti) suunnistettu}, jos $(e_1|_p,\dots,e_n|p)$ on (positiivisesti) suunnattu $T_pM$:n kanta kaikilla $p\in U$.

    Pisteittäinen suunnistus on \emph{jatkuva}, jos jokainen $p\in M$ on jossakin suunnistetussa kehyksessä.

    Moniston $M$ \emph{suunnistus} on jatkuva pisteittäinen suunnistus.
\end{maaritelma}